% ============================================================================
%  back/grammar.tex — درسنامهٔ جامعِ قواعد + مرورِ سریع
% ============================================================================
\Jpart{درسنامهٔ جامعِ قواعد و مرورِ سریع}
  {مُرَاجَعَةٌ شَامِلَةٌ لِلْقَوَاعِدِ}
  {همهٔ قواعدِ پایهٔ دوازدهم در یک نگاه — مناسبِ مرورِ شبِ امتحان}

% ============================================================================
%  ۱) حروف مشبهه فعل
% ============================================================================
\Jsec{١ \,\textbar\, اَلْحُرُوفُ الْمُشَبَّهَةُ بِالْفِعْلِ}

\begin{JQaede}{قاعدهٔ کلی}
حروفِ مشبههٔ فعل بر سرِ \textbf{جملهٔ اسمیّه} می‌آیند و آن را تأکید، پیوند،
تشبیه، استثناء، آرزو یا امید می‌بخشند. هنگامی‌که بر اسمِ ظاهر وارد می‌شوند:
\begin{itemize}
  \item \textbf{اسمِ} آن‌ها \textbf{منصوب} می‌شود.
  \item \textbf{خبرِ} آن‌ها \textbf{مرفوع} می‌ماند.
\end{itemize}
ترکیبِ «\textbf{لِأَنَّ}» به معنای «\textbf{زیرا، برای این‌که}» است و در
پاسخِ پرسش‌هایِ «چرا» می‌آید.
\end{JQaede}

\begin{center}
\renewcommand{\arraystretch}{1.6}
\begin{tabular}{@{}>{\JUIFont\bfseries\color{JPrimary}}l
                  >{\color{JInk}}p{0.30\linewidth}
                  >{\Ar\color{JPrimary!85!black}}p{0.34\linewidth}@{}}
\rowcolor{JPrimaryL}
\multicolumn{1}{@{}l}{\JUIFont\bfseries\small\color{JPrimary}الحَرْفُ} &
\JUIFont\bfseries\small\color{JPrimary}المَعْنَى &
\JUIFont\bfseries\small\color{JPrimary}مثال\\
إِنَّ & قطعاً، همانا، بی‌گمان (تأکید) & إِنَّ اللّٰهَ لا يُضِيعُ أَجْرَ
الْمُحْسِنِينَ\\
أَنَّ & که (پیوندِ دو جمله) & أَعْلَمُ أَنَّ اللّٰهَ عَلَىٰ كُلِّ شَيْءٍ
قَدِيرٌ\\
كَأَنَّ & گویی، مانندِ (تشبیه) & كَأَنَّهُنَّ الْيَاقُوتُ وَ الْمَرْجَانُ\\
لٰكِنَّ & ولی، اما (برطرف‌کردنِ ابهام) & إِنَّ اللّٰهَ لَذُو فَضْلٍ ...
وَلٰكِنَّ أَكْثَرَ النَّاسِ لا يَشْكُرُونَ\\
لَيْتَ & کاش (آرزو) & يَا لَيْتَنِي كُنْتُ تُرَابًا\\
لَعَلَّ & شاید، امید است & لَعَلَّكُمْ تَعْقِلُونَ\\
\end{tabular}
\end{center}

\begin{JDam}[دامِ آزمون]
پیش از ورودِ حرف: «الْعِلْمُ نَافِعٌ» (مبتدا و خبر، هر دو مرفوع)؛ پس از ورود:
«إِنَّ \textbf{الْعِلْمَ} \textbf{نَافِعٌ}» (اسمِ منصوب، خبرِ مرفوع). گزینه‌های
«الْعِلْمُ نَافِعًا» یا «الْعِلْمَ نَافِعًا» دام‌اند. همچنین در ترجمه به سیاق
توجّه کنید؛ «أَنَّ» معنای «که» دارد و دو جمله را پیوند می‌دهد.
\end{JDam}

% ============================================================================
%  ۲) لا نافیه جنس
% ============================================================================
\Jsec{٢ \,\text| لا النَّافِيَةُ لِلْجِنْسِ}

\begin{JQaede}{قاعدهٔ کامل}
معنای چهارمِ «لا»، «\textbf{هیچ... نیست}» است و بر سرِ \textbf{اسم} وارد
می‌شود:
\begin{itemize}
  \item \textbf{اسمِ لا:} معمولاً \textbf{بدونِ «ال»} و دارایِ \textbf{فتحه}
        است: لا \textbf{عِلْمَ} لَنَا...
  \item \textbf{خبرِ لا:} اگر اسمِ ظاهر باشد، \textbf{مرفوع} می‌ماند و اغلب
        \textbf{جارّ و مجرور} است: لا رَجُلَ \textbf{فِي الْحَفْلَةِ}.
  \item گاهی خبر \textbf{حذف} می‌شود: لا إِلٰهَ [مَوْجُودٌ] إِلَّا اللّٰهُ ،
        لا شَكَّ ، لا بَأْسَ.
\end{itemize}
معناهای دیگرِ «لا» برای مقایسه: «نه» در پاسخِ پرسش؛ لای نفیِ مضارع (نمی‌رود)؛
لای نهی (نرو) و «نباید» بر سرِ مضارعِ اوّل و سومِ شخص (لا يَذْهَبْ).
\end{JQaede}

\begin{JHadith}
\Ar لا شَيْءَ أَحْسَنُ مِنَ الْعَفْوِ عِنْدَ الْقُدْرَةِ.\quad
\Ar لا فَقْرَ كَالْجَهْلِ.\quad
\Ar لا رَيْبَ فِيهِ.

هیچ‌چیز زیباتر از بخشش هنگامِ توانایی نیست. هیچ فقری مانندِ نادانی نیست. هیچ
تردیدی در آن نیست.
\end{JHadith}

% ============================================================================
%  ۳) الحال
% ============================================================================
\Jsec{٣ \,\text| اَلْحَال (قیدِ حالت)}

\begin{JQaede}{قاعدهٔ کامل}
«حال» واژه‌ای است که \textbf{حالتِ یک اسم را هنگامِ وقوعِ فعل} نشان می‌دهد:
\begin{itemize}
  \item حال همیشه \textbf{منصوب} است.
  \item حال \textbf{نکره} است؛ اگر بعد از اسمِ معرفه، نکره‌ای مشابه بیاید،
        صفت نیست، بلکه حال است: رَأَيْتُ الْوَلَدَ، \textbf{مَسْرُورًا}.
  \item صفت: رَأَيْتُ الْوَلَدَ \textbf{الْمَسْرُورَ} (با «ال»، مطابقِ موصوف).
  \item ترجمه: با «در حالی‌که»، «به‌صورتِ»، «با حالتِ» یا قیدِ حالیِ فارسی.
  \item حالِ جمله‌ای: جملهٔ اسمیّه با \textbf{واوِ حالیّه} + ضمیر؛ در ترجمه
        اغلب ماضیِ استمراری: رَأَيْتُ الْفَلَّاحَ \textbf{وَ هُوَ يَجْمَعُ
        الْمَحْصُولَ}. (در حالی‌که محصول را جمع \emph{می‌کرد})
\end{itemize}
\end{JQaede}

% ============================================================================
%  ۴) مرورِ سریع (جدولِ یک‌نگاهی)
% ============================================================================
\Jsec{٤ \,\text| مُرَاجَعَةٌ سَرِيعَةٌ \,\text| مرورِ سریعِ قواعد}

\begin{center}
\renewcommand{\arraystretch}{1.65}
\begin{tabular}{@{}>{\JUIFont\bfseries\color{JPrimary}}p{0.20\linewidth}
                  >{\color{JInk}}p{0.38\linewidth}
                  >{\color{JInk}}p{0.36\linewidth}@{}}
\rowcolor{JPrimaryL}
\JUIFont\bfseries\small\color{JPrimary}موضوع &
\JUIFont\bfseries\small\color{JPrimary}قاعدهٔ کلیدی &
\JUIFont\bfseries\small\color{JPrimary}نمونهٔ کلیدی\\
حروفِ مشبهه & اسم: منصوب ؛ خبر: مرفوع &
إِنَّ الْعِلْمَ نَافِعٌ\\
لِأَنَّ & زیرا ، برای این‌که & لِأَنَّ الطَّرِيقَ مَسْدُودٌ\\
لا نافیهٔ جنس & اسم: بدونِ «ال» + فتحه ؛ خبر: مرفوع/جارّ و مجرور ؛ معنا:
«هیچ... نیست» & لا رَجُلَ فِي الْحَفْلَةِ\\
حال & منصوب + نکره ؛ ترجمه: «در حالی‌که» & خَرَجَ مَسْرُورًا\\
واوِ حالیّه & جملهٔ اسمیّه + واو + ضمیر & خَرَجَ وَ هُوَ يَبْتَسِمُ\\
\end{tabular}
\end{center}

% ============================================================================
%  ۵) قواعدِ درس‌های ۳ تا ۵ (الگو)
% ============================================================================
\Jsec{٥ \,\text| قواعدِ درسِ ٣ تا ٥ \,\text| محلِ تکمیل}

\begin{JQaede}[JCompact]{قواعدِ درسِ سوم (ثَلَاثُ قِصَصٍ قَصِيرَةٍ)}
[قواعدِ درسِ سوم اینجا درج شود — خلاصهٔ قاعده، جدول‌ها و مثال‌های کتاب.]
\end{JQaede}
\begin{JQaede}[JCompact]{قواعدِ درسِ چهارم (نِظَامُ الطَّبِيعَةِ)}
[قواعدِ درسِ چهارم اینجا درج شود.]
\end{JQaede}
\begin{JQaede}[JCompact]{قواعدِ درسِ پنجم (يَا إِلٰهِي)}
[قواعدِ درسِ پنجم اینجا درج شود.]
\end{JQaede}
